\chapter{Derivations}
\label{app:derivations}

This appendix gathers, in compact but complete form, derivations that the main text states or sketches.

\section{Solution of the capacitor equation}
\label{app:rc}

Equation~\eqref{eq:ode}: $C\dot V=(\Vs-V)/\Rs-V/\Rl-I_L$. Collect the terms in $V$:
\[
 C\dot V+\Big(\frac1{\Rs}+\frac1{\Rl}\Big)V=\frac{\Vs}{\Rs}-I_L
 \quad\Longrightarrow\quad
 \dot V+\frac V\tau=\frac{V_\infty}{\tau},
\]
with $\tau=C/(1/\Rs+1/\Rl)=(\Rs\parallel\Rl)C$ and $V_\infty=\tau(\Vs/\Rs-I_L)/C=(\Vs-I_L\Rs)\,\Rl/(\Rl+\Rs)$. Multiplying by the integrating factor $\e^{t/\tau}$:
\[
 \frac{\mathrm d}{\mathrm dt}\big(V\e^{t/\tau}\big)=\frac{V_\infty}{\tau}\e^{t/\tau}
 \ \Rightarrow\
 V(t)\e^{t/\tau}=V_0+V_\infty(\e^{t/\tau}-1)
 \ \Rightarrow\
 V(t)=V_\infty+(V_0-V_\infty)\e^{-t/\tau}.
\]
Setting $V(t_{\mathrm{th}})=V_{\mathrm{th}}$ and solving gives Eq.~\eqref{eq:tth}.

\paragraph{Discrete step.} Explicit Euler on the source term alone gives $V_{k+1}=V_k+\frac{\Delta t}{\Rs C}(\Vs-V_k)$, with amplification factor $1-\Delta t/(\Rs C)\in(0,1)$ when $\Delta t<\Rs C$ (stable, here $\Delta t/\Rs C=6.7\times10^{-4}$). Leakage alone is $V\to V\e^{-\Delta t/\Rl C}$ exactly. Composing them (Eq.~\eqref{eq:chargestep}) commits an $O(\Delta t^2)$ local error per step, hence $O(\Delta t)$ globally, with a coefficient of order $\Delta t/\tau\sim10^{-3}$: we measured an error of $1.7\times10^{-4}$ V after $\SI{1}{s}$.

\section{Maximum power into the capacitor}
Power into the capacitor through $\Rs$ at voltage $V$ is $p(V)=V(\Vs-V)/\Rs$. Then $p'(V)=(\Vs-2V)/\Rs=0$ at $V=\Vs/2$ and $p_{\max}=\Vs^2/(4\Rs)=P_h/4$, with $P_h=\Vs^2/\Rs$. Equivalently $p=P_h\,x(1-x)$, $x=V/\Vs$.

\section{Duty-cycle theorem}
See Section~\ref{sec:hysteresis}. With $T_c=\Delta E/P_{\mathrm{in}}$, $T_a=\Delta E/(P_{\mathrm{act}}-P_{\mathrm{in}})$: $D=T_a/(T_a+T_c)=P_{\mathrm{in}}/P_{\mathrm{act}}$, because
$\dfrac1{P_{\mathrm{act}}-P_{\mathrm{in}}}\Big/\Big(\dfrac1{P_{\mathrm{act}}-P_{\mathrm{in}}}+\dfrac1{P_{\mathrm{in}}}\Big)=\dfrac{P_{\mathrm{in}}}{P_{\mathrm{in}}+(P_{\mathrm{act}}-P_{\mathrm{in}})}$.

\section{Framed slotted ALOHA: slot statistics}
\label{app:fsa}

Each of $n$ nodes picks one of $S$ slots independently and uniformly. The occupancy $X$ of a given slot is Binomial$(n,1/S)$:
$\Pr[X=k]=\binom nk S^{-k}(1-1/S)^{n-k}$. Hence
$\Pr[X=0]=(1-1/S)^n$, $\Pr[X=1]=\frac nS(1-1/S)^{n-1}$. By linearity of expectation over the $S$ slots, $E_0=S\Pr[X=0]$, $E_1=S\Pr[X=1]=n(1-1/S)^{n-1}$ and $E_c=S-E_0-E_1$. A tagged node succeeds iff each of the other $n-1$ nodes avoided its slot: $(1-1/S)^{n-1}$.

\paragraph{Optimal frame size.} Let $r=n/S$. For large $S$, $(1-1/S)^{n-1}\to\e^{-r}$, so $\eta_S\to r\e^{-r}$, with $\mathrm d\eta_S/\mathrm dr=(1-r)\e^{-r}=0$ at $r=1$ and $\eta_S(1)=1/\e$.

\paragraph{Pure and slotted ALOHA.} A transmission starting at $t_0$ is vulnerable to any other start in $(t_0-T,t_0+T)$ (pure) or in the same slot (slotted). With Poisson arrivals of rate $g$ per packet time, the probability of no other start in an interval of length $2T$ (resp.\ $T$) is $\e^{-2g}$ (resp.\ $\e^{-g}$). The maximum of $g\e^{-2g}$ is at $g=1/2$ ($1/2\e$) and of $g\e^{-g}$ at $g=1$ ($1/\e$).

\section{SIC: the ratio rule}
\label{app:sicratio}

For perfect cancellation and negligible noise, powers $P_{(k)}=P_{(1)}\rho^{-(k-1)}$ with $\rho>1$. At stage $k$ the interference is the sum of all weaker packets: $P_{(k)}\sum_{j\ge1}\rho^{-j}$ for an infinite chain, which converges to $P_{(k)}/(\rho-1)$. Thus $\gamma_{(k)}=\rho-1$ for every $k$, and decoding all packets requires $\rho-1\ge\gamma_{\mathrm{req}}$, i.e.\ $\rho\ge1+\gamma_{\mathrm{req}}$ (in linear units). For a finite chain of $m$ packets the ratio at the last stages is larger, so $\rho\ge1+\gamma_{\mathrm{req}}$ is sufficient. With $\gamma_{\mathrm{req}}=3$~dB ($\approx2.0$) we need $\rho\ge3.0$ ($4.8$~dB). The strongest of four packets exceeds the weakest by $\rho^{3}=27$ ($14.3$~dB).

\section{Sidebands of a switched reflection}
\label{app:fourier}

Let $s(t)\in\{0,1\}$ be a unit square wave with period $1/f_s$ and 50\% duty. Its Fourier series is $s(t)=\tfrac12+\tfrac2\pi\sum_{k=1,3,5,\dots}\tfrac1k\sin(2\pi kf_st)$. The reflection coefficient is $\Gamma(t)=\Gamma_0+\Delta\Gamma\,s(t)$ and the reflected wave is $y(t)=\Gamma(t)\,A\cos(2\pi f_ct)$. Using $\sin a\cos b=\tfrac12[\sin(a+b)+\sin(a-b)]$ on the fundamental ($k=1$) term:
\[
 \Delta\Gamma\frac{2}{\pi}A\sin(2\pi f_st)\cos(2\pi f_ct)=\frac{\Delta\Gamma}{\pi}A\big[\sin 2\pi(f_c+f_s)t-\sin2\pi(f_c-f_s)t\big].
\]
Each sideband has amplitude $\Delta\Gamma A/\pi$, hence a power ratio of $(\Delta\Gamma/\pi)^2$ to the incident carrier ($A$). For $\Delta\Gamma=1$ this is $1/\pi^2=0.101$ ($-9.9$~dB); for $\Delta\Gamma=2$ it is $4/\pi^2=0.405$ ($-3.9$~dB). The constant term $\tfrac12\Delta\Gamma+\Gamma_0$ multiplies the carrier itself, giving the unshifted ``static'' reflection of Eq.~\eqref{eq:y0}.

\section{Friis in two lines}
Isotropic power density at range $d$: $S=P_tG_t/(4\pi d^2)$. Effective aperture of a receive antenna of gain $G_r$: $A_e=G_r\lambda^2/(4\pi)$ (a standard result obtained from antenna reciprocity and thermal equilibrium; see, e.g., \citet{pozar2011}). The received power is $S A_e=P_tG_tG_r(\lambda/4\pi d)^2$.

\section{3GPP UMa continuity at the breakpoint}
At $d_{2D}=d'_{BP}$ and ignoring the height difference in $d_{3D}$, $\mathrm{PL}_1=28+22\log d'_{BP}+20\log f_c$ and $\mathrm{PL}_2=28+40\log d'_{BP}+20\log f_c-9\log(d_{BP}'^{2}+\Delta h^2)\approx28+40\log d'_{BP}+20\log f_c-18\log d'_{BP}=28+22\log d'_{BP}+20\log f_c$ when $\Delta h\ll d'_{BP}$: the two branches agree, confirming the continuity claim. For the demo ($d'_{BP}=\SI{147.8}{m}$) we measured $75.17$ and $\SI{75.19}{dB}$ at $147.8$ and $\SI{148}{m}$.
